\section{Actual framed residue coefficients under arbitrary original scheme pullback}

\begin{lemma}\label{lem:fc-native-unit-section-morphism}
\lean{EnvelopeProperty.Berkovich.SchemeFramedSectionPullback.affinePullbackUnitTop_eq_sectionAtOne,EnvelopeProperty.Berkovich.SchemeFramedSectionPullback.affinePullbackUnitTop_naturality,EnvelopeProperty.Berkovich.SchemeFramedSectionPullback.affinePullbackUnitTop_comp,EnvelopeProperty.Berkovich.SchemeFramedSectionPullback.transportedSection,EnvelopeProperty.Berkovich.SchemeFramedSectionPullback.transportedSection_sectionHom}
\leanok
\uses{lem:fc-original-line-section-fibre,lem:fc-original-section-pullback-frame,scheme:existing-affine-module-pullback-sections}
For every actual scheme map $g:Y\to Z$, original module $M$ and
original whole section $\sigma\in\Gamma(Z,M)$, the native pullback-unit
section is exactly the value at one of the genuine pullback of its
original structure-line section morphism. For every actual
isomorphism $\alpha:g^*M\simeq N$, define the actual transported
whole section $\tau$ by applying $\alpha$ to that native unit section.
Its original section morphism is precisely the genuine pulled section
morphism followed by $\alpha$.
\end{lemma}
\begin{proof}
Naturality of the native pullback adjunction unit for the original
structure-line section morphism identifies its original unit image
with that pulled morphism applied to the image of one. The genuine
native structure-unit comparison identifies that image with one.
Transport from the inverse image of the top open to the top open
uses native presheaf naturality, so the actual native whole section
is exactly the asserted value at one. The original section and
structure-line-morphism equivalence is inverse to evaluation at one;
applying it to the actual image under $\alpha$ gives the final whole
section-morphism identity.
\end{proof}

\begin{definition}\label{def:fc-induced-native-residue-frame}
\lean{EnvelopeProperty.Berkovich.SchemeFramedSectionPullback.inducedResidueFrame,EnvelopeProperty.Berkovich.SchemeFramedSectionPullback.inducedResidueFrame_compat,EnvelopeProperty.Berkovich.SchemeFramedSectionPullback.inducedResidueFrame_eq,EnvelopeProperty.Berkovich.SchemeFramedSectionPullback.inducedResidueFrame_hom}
\leanok
\uses{berkovich:actual-residue-module-pullback-iso,berkovich:actual-residue-module-pullback-frame}
Fix any actual native valued point $y$ of $Y$ and any actual original
residue frame $e_M$ for $M$ at $g(y)$. Pull $e_M$ through the actual
residue-field map $\rho:\kappa_Z(g(y))\to\kappa_Y(y)$ and the genuine
native structure-unit comparison. Precede this by the inverses of
the actual residue-module pullback comparison and the actual
residue pullback of $\alpha$. The result is an entire actual frame
$e_N$ of the same original $N$-fibre at $y$. The whole composition
of the residue-module comparison, pulled $\alpha$, and $e_N$ is
exactly the actual residue-field pullback of $e_M$ with its genuine
native structure-unit comparison.
\end{definition}
\begin{proof}
All four maps are actual native module isomorphisms on the original
residue-field spectrum, so their composite is the stated frame.
Compose it with the original residue-module comparison and the
pulled $\alpha$. The adjacent original isomorphisms and their
inverses cancel by the native inverse identities. The remaining
whole map is exactly the actual residue pullback of $e_M$ followed
by the native unit comparison. No frame-preservation equality is
assumed.
\end{proof}

\begin{theorem}\label{thm:fc-native-framed-section-residue-transport}
\lean{EnvelopeProperty.Berkovich.SchemeFramedSectionPullback.residueFrameCoefficient,EnvelopeProperty.Berkovich.SchemeFramedSectionPullback.transportedSection_coefficient}
\leanok
\uses{lem:fc-native-unit-section-morphism,def:fc-induced-native-residue-frame,lem:fc-original-line-section-frame-naturality,berkovich:actual-residue-section-hom-pullback}
For every original whole section $\sigma$, form the actual transported
section $\tau$ just defined. Pull these original sections to their
actual native residue fibres and apply their actual frames. If
$c_M\in\kappa_Z(g(y))$ and $c_N\in\kappa_Y(y)$ are the resulting
original coefficients, then
\[
 c_N=\rho(c_M).
\]
The identity uses the derived entire original frame $e_N$ and holds
for every original section, including zero. It assumes no section,
coefficient, frame-preservation or residue-evaluation certificate.
\end{theorem}
\begin{proof}
Use the first lemma to express each native residue-unit section as
the value at one of its genuine pulled section morphism. Native
pullback composition and the actual residue-field square identify
the two routes for pulling the original structure-line morphism.
After the actual pulled $\alpha$ is inserted, composing with the
derived frame cancels $\alpha$ and the residue-module comparison,
by the preceding whole-frame identity. The resulting whole map is
the genuine pullback through $\rho$ of the original framed scalar
endomorphism. Every original structure-line endomorphism is
multiplication by its actual value on one. Native scalar pullback
therefore gives multiplication by the actual pulled scalar.
Evaluate this whole original identity at one and use naturality of
the genuine affine section isomorphism for the actual residue-field
map. This gives precisely $c_N=\rho(c_M)$.
\end{proof}

\begin{corollary}\label{cor:fc-native-framed-section-valuation-transport}
\lean{EnvelopeProperty.Berkovich.SchemeFramedSectionPullback.transportedSection_valuation}
\leanok
\uses{thm:fc-native-framed-section-residue-transport,scheme:valued-points}
The actual native valuation of $c_N$ at $y$ equals the actual native
valuation of $c_M$ at the whole original mapped valued point $g(y)$.
No properness, compactness, bounded base valuation or model norm is
needed, and zero coefficients are included.
\end{corollary}
\begin{proof}
Apply the original residue valuation at $y$ to the preceding exact
coefficient identity. By definition the residue valuation of the
whole native mapped point $g(y)$ is its original residue valuation
composed with the actual residue-field map $\rho$. This is exactly
the claimed equality, also at zero.
\end{proof}
